% Part: normal-modal-logic
% Chapter: filtrations
% Section: S5-fmp

\documentclass[../../../include/open-logic-section]{subfiles}

\begin{document}

\olfileid[pt]{nml}{fil}{fmp}

\olsection{\Log{K} e \Log{S5} têm a propriedade do modelo finito}

\begin{defn}
  Diz-se que um sistema $\Sigma$ de lógica modal tem a \emph{propriedade
    do modelo finito} se, sempre que !!a{formula}~$!A$ é verdadeira em um
  mundo de um modelo de $\Sigma$, $!A$ é verdadeira em um mundo de algum
  modelo \emph{finito} de~$\Sigma$.
\end{defn}

\begin{prop}\ollabel{prop:K-fmp}
  \Log{K} tem a propriedade do modelo finito.
\end{prop}

\begin{proof}
  \Log{K} é o conjunto das !!{formula}s válidas, O.s.\@, todo modelo é
  um modelo de~\Log{K}. Por \olref[fil]{thm:filtrations}, se
  $\mSat{M}{!A}[w]$, então $\mSat{M^*}{!A}[w]$ para toda filtração
  de~$\mModel{M}$ por meio do conjunto $\Gamma$ das sub-!!{formula}s
  de~$!A$. Toda !!{formula} tem apenas um número finito de
  sub-!!{formula}s; logo, $\Gamma$ é finito. Por
  \olref[fin]{prop:filt-are-finite}, $\card{W^*} \le 2^n$, em que $n$
  é o número de !!{formula}s em~$\Gamma$. E, como \Log{K} não impõe
  restrição alguma aos modelos, $\mModel{M^*}$ é um modelo de \Log{K}.
\end{proof}

Para mostrar por meio de filtrações que uma lógica~\Log{L} tem a
propriedade do modelo finito, é essencial que a filtração de um modelo
de \Log{L} seja ela própria um modelo de \Log{L}. Isso frequentemente
exige bastante trabalho, e nem toda filtração produz um modelo de
\Log{L}. Contudo, para modelos universais, isso ainda vale.

\begin{prop}\ollabel{prop:univ-fin}
  Seja $\mClass{U}$ a classe dos modelos universais (veja
  \olref[frd][es5]{prop:S5=univ}) e seja $\mClass{U}_\mathrm{Fin}$ a
  classe de todos os modelos universais finitos. Então, toda
  !!{formula}~$!A$ é válida em $\mClass{U}$ se, e somente se, é válida em
  $\mClass{U}_\mathrm{Fin}$.
\end{prop}

\begin{proof}
  Modelos universais finitos são modelos universais; portanto, a direção
  da esquerda para a direita é trivial. Para a direção da direita para a
  esquerda, suponha que~$!A$ seja falsa em algum mundo $w$ de um modelo
  universal $\mModel{M}$. Seja $\Gamma$ um conjunto que contenha $!A$ e
  todas as suas subfórmulas; evidentemente, $\Gamma$ é finito. Tome uma
  filtração $\mModel{M^*}$ de $\mModel{M}$; então $\mModel{M^*}$ é finito
  por \olref[fin]{prop:filt-are-finite} e, por
  \olref[fil]{thm:filtrations}, $!A$ é falsa em $[w]$ em
  $\mModel{M^*}$. Resta observar que $\mModel{M^*}$ também é universal:
  dados $u$ e $v$, por hipótese $Ruv$ e, por
  \olref[fil]{defn:filtration}\olref[fil]{defn:filtration-R}, também
  $R^*[u][v]$.
\end{proof}

\begin{cor}\ollabel{cor:S5fmp}
  \Log{S5} tem a propriedade do modelo finito.
\end{cor}

\begin{proof}
  Por \olref[frd][es5]{prop:S5=univ}, se $!A$ é verdadeira em um mundo
  de algum modelo reflexivo e euclidiano, então é verdadeira em um mundo
  de um modelo universal. Por \olref{prop:univ-fin}, ela é verdadeira em
  um mundo de um modelo universal finito (a saber, a filtração do modelo
  por meio do conjunto das sub-!!{formula}s de~$!A$). Todo modelo
  universal também é reflexivo e euclidiano; logo, $!A$ é verdadeira em
  um mundo de um modelo reflexivo e euclidiano finito.
\end{proof}

\begin{prob}
  Mostre que toda filtração de um modelo serial ou reflexivo também é,
  respectivamente, serial ou reflexiva.
\end{prob}

\begin{prob}
  Encontre uma filtração não simétrica (não transitiva, não euclidiana)
  de um modelo simétrico (transitivo, euclidiano).
\end{prob}

\end{document}
